\subsection{Monomial Ideals and Dickon's Lemma}

\begin{definition}[title = Definition 2.3.1]
    For a (potentially infinite) family $f_s \ (s \in S)$ of polynomials indexed by $s \in S$, we denote $\langle f_s: s\in S \rangle$ the ideal generated by these polynomials, 
    i.e. the ideal 
    \[
      \left\{ \sum_{s\in S} f_sh_s: h \in k[x_1, \dots, x_n] \text{ and } h_s \neq 0 \text{ for only finitely many } s \in S \right\} 
    \]
\end{definition}

\begin{definition}[title = Definition 2.3.2]
    The \textit{monomial ideal} for a set $A \subseteq \mathbb{Z}^n_{\geq 0}$ of exponent vectors is defined as $\langle x^{\alpha}: \alpha \in A \rangle$.
\end{definition} 

\begin{definition}[title = Definition 2.3.3]
    In $\mathbb{Z}^n_{\geq 0}$ we introduce the partial order $\geq$ by $\alpha \geq \beta :\Leftrightarrow \alpha_1 \geq \beta_1, \ldots, \alpha_n \geq \beta_n$ 
    with $\alpha = (\alpha_1, \ldots, \alpha_n), \beta = (\beta_1, \ldots, \beta_n) \in \mathbb{Z}^n_{\geq 0}$.
\end{definition} 

\begin{example}[title = Example]
\begin{minipage}{0.49\textwidth}
    \begin{tikzpicture}
  % Grid
  \draw[step=1, gray!40, very thin] (0,0) grid (5,5);

  % Achsen
  \draw[->, thick] (0,0) -- (5.4,0) node[right] {$x$};
  \draw[->, thick] (0,0) -- (0,5.4) node[above] {$y$};

  % Beschriftung der Achsen
  \foreach \i in {0,...,5} {
    \draw (\i,0.08) -- (\i,-0.08) node[below] {\i};
    \draw (0.08,\i) -- (-0.08,\i) node[left] {\i};
  }

  % Punkt (1,1)
  \fill[kGreen] (0,1) circle (2.5pt);
  \fill[kGreen] (2,2) circle (2.5pt);
  \fill[kGreen] (3,4) circle (2.5pt);
\end{tikzpicture} 

Partial order: $(0,1) \leq (2,2) \leq (3,4)$
\end{minipage}
\hspace*{0.2cm}
\begin{minipage}{0.49\textwidth}
    \begin{tikzpicture}
  % Grid
  \draw[step=1, gray!40, very thin] (0,0) grid (5,5);

  % Achsen
  \draw[->, thick] (0,0) -- (5.4,0) node[right] {$x$};
  \draw[->, thick] (0,0) -- (0,5.4) node[above] {$y$};

  % Beschriftung der Achsen
  \foreach \i in {0,...,5} {
    \draw (\i,0.08) -- (\i,-0.08) node[below] {\i};
    \draw (0.08,\i) -- (-0.08,\i) node[left] {\i};
  }

  % Punkt (1,1)
  \fill[kGreen] (0,1) circle (2.5pt);
  \fill[kGreen] (2,2) circle (2.5pt);
  \fill[kGreen] (3,4) circle (2.5pt);
  \fill[kGreen] (3,1) circle (2.5pt);
\end{tikzpicture} 

No Partial order: $(2,2) \nleq (3,1)$ and $(3,1) \nleq (2,2)$
\end{minipage}
\end{example} 

\begin{definition}[title = Definition 2.3.4]
    The \textit{Minkowski Sum} $A + B$ of $A, B \in \mathbb{Z}^n_{\geq 0}$ is defined as $A + B := \{\alpha + \beta: \alpha \in A, \beta \in B\}$. 
\end{definition}

\begin{example}[title = Example]
\begin{minipage}[t]{0.49\textwidth}
    \begin{tikzpicture}

  % Achsen
  \draw[->, thick] (0,0) -- (2.4,0) node[right] {$x$};
  \draw[->, thick] (0,0) -- (0,2.4) node[above] {$y$};

  % Achsenbeschriftung
  \foreach \i in {0,...,2} {
    \draw (\i,0.08) -- (\i,-0.08) node[below] {\i};
    \draw (0.08,\i) -- (-0.08,\i) node[left] {\i};
  }

  % Dreieck
  \draw[thick, red] (0,0) -- (1,0) -- (0,1) -- cycle;

  % Punkte
  \fill[kGreen] (0,0) circle (2.5pt);
  \fill[kGreen] (1,0) circle (2.5pt);
  \fill[kGreen] (0,1) circle (2.5pt);
\end{tikzpicture}

\[A = \{(0,0), (1,0), (0,1)\}\]
\end{minipage}
\hspace*{0.2cm}
\begin{minipage}[t]{0.49\textwidth}
\begin{tikzpicture}

  % Achsen
  \draw[->, thick] (0,0) -- (5.4,0) node[right] {$x$};
  \draw[->, thick] (0,0) -- (0,2.4) node[above] {$y$};

  % Achsenbeschriftung
  \foreach \i in {0,...,5} {
    \draw (\i,0.08) -- (\i,-0.08) node[below] {\i};
  }
  \foreach \i in {1,2} {
    \draw (0.08,\i) -- (-0.08,\i) node[left] {\i};
  }

  % Dreieck
  \draw[thick, blue] (0,0) -- (1,1) -- (4,0) -- cycle;

  % Punkte
  \fill[kGreen] (0,0) circle (2.5pt);
  \fill[kGreen] (1,1) circle (2.5pt);
  \fill[kGreen] (4,0) circle (2.5pt);
\end{tikzpicture}
\[B = \{(0,0), (1,1), (4,0)\}\]
\end{minipage}

\begin{center}
    \begin{tikzpicture}

  % Achsen
  \draw[->, thick] (0,0) -- (5.4,0) node[right] {$x$};
  \draw[->, thick] (0,0) -- (0,2.4) node[above] {$y$};

  % Achsenbeschriftung
  \foreach \i in {0,...,5} {
    \draw (\i,0.08) -- (\i,-0.08) node[below] {\i};
  }
  \foreach \i in {1,2} {
    \draw (0.08,\i) -- (-0.08,\i) node[left] {\i};
  }

  % Dreieck
  \draw[thick, blue] (0,0) -- (1,1) -- (4,0) -- cycle;
  \draw[thick, red] (0,0) -- (1,0) -- (0,1) -- cycle;
  \draw[thick, red] (1,1) -- (2,1) -- (1,2) -- cycle;
  \draw[thick, red] (4,0) -- (4,1) -- (5,0) -- cycle;

  % Punkte
  \fill[kGreen] (0,0) circle (2.5pt);
  \fill[kGreen] (1,1) circle (2.5pt);
  \fill[kGreen] (4,0) circle (2.5pt);

  \fill[kGreen] (4,0) circle (2.5pt);
  \fill[kGreen] (4,1) circle (2.5pt);
  \fill[kGreen] (5,0) circle (2.5pt);

  \fill[kGreen] (1,1) circle (2.5pt);
  \fill[kGreen] (2,1) circle (2.5pt);
  \fill[kGreen] (1,2) circle (2.5pt);

  \fill[kGreen] (0,0) circle (2.5pt);
  \fill[kGreen] (1,0) circle (2.5pt);
  \fill[kGreen] (0,1) circle (2.5pt);
\end{tikzpicture}
\[A + B = \{(0,0), (1,0), (0,1), (1,1), (2,1), (1,2), (4,0), (4,1), (5,0)\}\]
\end{center}
\end{example} 


\begin{lemma}[title = Lemma 2.3.5]
    For $A \subseteq \mathbb{Z}^n_{\geq 0}$, we have $\langle x^{\alpha}: \alpha \in A \rangle = lin_k \{x^{\beta}: \beta \in A + \mathbb{Z}^n_{\geq 0}\}$ with $x^{\beta} = \{\sum_{\beta}c_{\beta} \}$ %todo(foto)  
    where $lin_k$ is the linear hul with respect to the field $k$.  
    
    \paragraph{Proof:} Let's verify \enquote{$\geq$}: IF $\beta \in A + \mathbb{Z}^n_{\geq 0}$, then $\beta  = \alpha + \gamma$ with $\gamma \in \mathbb{Z}^n_{\geq 0}$. 
    But then 
    \[x^{\beta} = x^{\alpha + \gamma} = \underbrace{\underbrace{x^{\alpha}}_{\mathclap{\text{in the ideal of the l.h.s}}} \cdot x^{\gamma}}_{\mathclap{\text{in the ideal on the l.h.s}}}\]

    $\Rightarrow$ every $x^{\beta}$ with $\beta \in A + \mathbb{Z}^n_{\geq 0}$ is in the l.h.s ideal \\
    $\Rightarrow \quad lin_k \{x^{\beta}: \beta \in A + \mathbb{Z}^n_{\geq 0}\}$ is in the l.h.s ideal \\[0.5em]
    Now let's verify $\leq$: Consider a product of any $x^{\alpha}, (\alpha \in A)$ with any $h = \sum_{\gamma} c_{\gamma} x^{\gamma} \in k[x_1, \dots, x_n]$.

    \[x^{\alpha} \cdot h = \sum_{\gamma} c_{\gamma} \cdot x^{\alpha} \cdot x^{\gamma} = \sum_{\gamma} c_{\gamma} x^{\alpha + \gamma} \in lin_k \{x^{\beta}: \beta \in A + \mathbb{Z}^n_{\geq 0}\}\]

    $\Rightarrow$ all such products belong to the r.h.s \\
    $\Rightarrow$ their linear combinations also belong to the r.h.s \\
      \begin{flushright}
        $\Box $
      \end{flushright}
\end{lemma}

\begin{example}[title = Example]  
    \begin{center}
    \begin{tikzpicture}
        \fill[red!10] (1,5) -- (5,5) -- (5,1) -- (3,1) -- (3,3) -- (2,3) -- (2,4) -- (1,4) -- cycle;
        % Grid
        \draw[step=1, gray!40, very thin] (0,0) grid (5,5);
        
        % Achsen
        \draw[->, thick] (0,0) -- (5.4,0) node[right] {$x_1$};
        \draw[->, thick] (0,0) -- (0,5.4) node[above] {$x_2$};
        
        % Beschriftung der Achsen
        \foreach \i in {0,...,5} {
            \draw (\i,0.08) -- (\i,-0.08) node[below] {\i};
            \draw (0.08,\i) -- (-0.08,\i) node[left] {\i};
            }
            
            
            %ideal
            \fill[red] (1,4) circle (2.5pt);
            \fill[red] (2,3) circle (2.5pt);
            \fill[red] (3,1) circle (2.5pt); 
            
            %polynomial 1
            \fill[kGreen] (2,4) circle (2.5pt);
            \fill[kGreen] (4,5) circle (2.5pt);
            \fill[kGreen] (3,2) circle (2.5pt);  
            
            %polynomial 2
            \fill[blue] (4,1) circle (2.5pt);
            \fill[blue] (2,2) circle (2.5pt);
        \end{tikzpicture} 
    \end{center}
    \begin{minipage}[t]{0.49\textwidth}
        $A = \{(1,4), (2,3), (3,1)\}$ \\
        and the ideal:
        $\textcolor{red}{I = \langle x_1x_2^4, x_1^2x_2^3, x_1^3x_2 \rangle}$\\[0.3dem]
        and the Minkowski sum: $A + \mathbb{Z}^2_{\geq 0}$ (marked in red)
    \end{minipage}
    \hspace*{1cm}
    \begin{minipage}[t]{0.39\textwidth}  
    $\textcolor{kGreen}{p_1 = 3x_1^3x_2^2 + 7x_1^2x_2^4 - 6x_1^4x_2^5}$ \\[0.3em]
    Is tis polynomial in the ideal? \\
    $\rightarrow$ Yes! \\[0.5em]
    $\textcolor{blue}{p_2 = x_1^2x_2^2-x_1^4x_2^1}$ \\
    Is this polynomial in the ideal? \\
    $\rightarrow$ No! 
    \end{minipage}

\end{example}